\section{Related Work}
\label{sec:related}

The push/pause/hold protocol draws on active-probing paradigms from
several fields where injecting a known signal into a medium and
observing the response is the standard approach to discovering
hidden structure.

\paragraph{Radar and sonar.}
Constant-false-alarm-rate detection was developed for radar signal
processing~\citep{finn1968cfar}. A transmitter emits a known waveform;
a receiver listens for echoes. The CFAR threshold adapts to local
noise conditions, maintaining a constant false-alarm probability
regardless of background noise level. Active sonar applies the same
principle underwater: a sound pulse propagates through the medium, and
the return signal reveals structure (seafloor, submarines, fish
schools) that passive listening cannot detect. Our protocol applies
this paradigm to coupled systems: the push is the transmitted pulse,
the receiver's objective is the return signal, and CFAR discriminates
the coupling step from intrinsic noise.

\paragraph{Seismology.}
Controlled-source seismology detonates explosives or vibrates the
ground at known locations, then measures ground motion at distant
sensors. The subsurface structure---rock layers, fault lines, fluid
reservoirs---is inferred from how the seismic waves propagate and
reflect. The structure cannot be observed passively; it must be
excited. Our method follows the same logic: coupling structure is
inferred from how a perturbation propagates from one agent to
another through the substrate.

\paragraph{Functional brain imaging (fMRI).}
In functional MRI, researchers study brain activity by presenting
controlled stimuli and measuring the hemodynamic response---the
slow, sustained blood-oxygen change that follows neural activation.
The hemodynamic response function (HRF) is a low-pass filter that
smears brief neural events into sustained signal changes. This led
to the \emph{block design} paradigm~\citep{friston1999fmri}:
sustained stimulation blocks of 20--30 seconds alternating with rest
periods, rather than single brief events. Our push/pause blocks are
the same design: a sustained perturbation block that survives the
substrate's thermal mass or response latency, followed by a pause
block that captures recovery. The ABA structure (stimulus, rest,
stimulus) is standard practice in experimental psychology and
neuroscience~\citep{box1978statistics}.

\paragraph{Medical percussion.}
The oldest active-probing diagnostic: a physician taps the body and
listens to the resonance. Solid tissue, fluid-filled cavities, and
air-filled spaces produce different sounds. The technique dates to
the 18th century~\citep{auerbrugger1761percussion}. The principle---
perturb the system, listen to the response, infer internal
structure---is exactly ours, applied to coupling channels between
autonomous agents.

\paragraph{Causal discovery.}
Pearl's do-calculus~\citep{pearl2009causality} provides the
theoretical framework for causal inference from interventions.
Convergent cross-mapping~\citep{sugihara2012ccm} infers coupling
from time series but fails when systems are actively exploring.
Neural relational inference~\citep{kipf2018nri} learns interaction
graphs from trajectories but operates passively. Our method differs
in that the interventions are performed by autonomous agents on
themselves, coordinated through a shared lease, rather than by a
central researcher.

\paragraph{Network tomography.}
Network tomography~\citep{vardi1996tomography} infers internal
network properties from end-to-end measurements using active probe
packets. Active topology inference~\citep{coates2002tomography}
selects measurement paths adaptively to reduce uncertainty about
network structure. These approaches share our active-probing
philosophy but assume a single centralized prober with an explicit
identification objective. Our agents are autonomous---each with its
own objective---and the probing is a cooperative side effect of
coordination, not the agents' primary goal.
